\documentclass[docs/main.tex]{subfiles}


\usepackage{parskip} % no indents

\newcommand{\reviewercomment}[1]{\vspace{0.3cm}\noindent\textbf{#1}\vspace{0.15cm}}
\newcommand{\manuscriptquote}[1]{\begin{quote}\textit{#1}\end{quote}}

\begin{document}

Dear Imaging Neuroscience Editors and Reviewers,

\vspace{0.5cm}

We thank you for the excellent and thoughtful reviews on our manuscript, ``Estimating fMRI Timescale Maps.'' We deeply appreciate the time and care the reviewers spent in providing such detailed feedback. We hope our appreciation is reflected in the careful revisions we have undertaken and are confident that these changes have made the work significantly clearer and more accessible, both for audiences focused on the statistical modeling and for practitioners applying these methods. We have revised the manuscript to address all the reviewers' comments. As an overview, we have:

\begin{enumerate}
    \item Incorporated \textit{RAPIDtide} denoising of the HCP dataset to mitigate physiological artifacts and systemic low-frequency oscillations, and regenerated all \textit{Data Analysis Results }(4.2; Reviewer 2).
    \item Added a new section, \textit{Practical Considerations for fMRI Timescale Estimation} (5), with three subsections explicitly addressing the effects of the \textit{hemodynamic response function} (5.1), \textit{measurement noise} (5.2), and \textit{periodicity} (5.3) on timescale estimation (Reviewers 1-3). 
    \item Included additional simulation settings to demonstrate the impact of measurement noise (Figure 7) and damped oscillations (Figure 8) on our proposed timescale methods (Reviewers 2-3).
    \item In the methods section (2), we prefaced reliance on Hansen’s textbook, revised the notation to better distinguish between the time- and autocorrelation-domain methods and their respective estimands and estimators, included more background on \textit{Estimation and Inference under Misspecification} (2.1.3), and introduced \textit{Estimator Properties} (2.4) before \textit{Standard Error of the Estimators} (2.5; Reviewer 3)
\end{enumerate}

\vspace{0.5cm}

Below, we provide a point-by-point response to the reviewers’ comments. Reviewer comments are in \textbf{bold}, our responses in plain text, and changes to the manuscript in \textit{italics}.

\subsection*{Reviewer 1}

\reviewercomment{Should the manuscript require any revision, I would like to suggest the following. Because the paper stands to become the de facto technical reference for timescale estimation with fMRI, it would be nice for the authors to also comment on fMRI-specific methodological considerations as well as their interaction with statistical procedures discussed here (please note that I am not requesting additional analyses).}

\reviewercomment{1. In the former case, for instance, the authors might simply mention potential artifacts that are worth considering in the context of timescale estimation, such as physiological confounds and head motion (Goldberg et al., 2024).}

We agree that physiological confounds and head motion are critical considerations. We have added a discussion on this in \textbf{Section 5.2 (Effects of Measurement Noise)} and cited Goldberg et al. (2024). Furthermore, we now explicitly mention using RAPIDtide to mitigate physiological artifacts in \textbf{Section 4.1}.

\manuscriptquote{To further mitigate systemic low-frequency physiological oscillations driven by respiration and CO2 variation, we applied RAPIDtide (Frederick, 2025). [Section 4.1]}

\manuscriptquote{Therefore in practice, effective denoising is a prerequisite for separating neural timescales from artifacts such as head motion or physiological noise. [Section 5.2]}

\reviewercomment{2. A comment on low-SNR regions is also advisable: the "limbic network" is here observed to have short autocorrelation times, and this is attributed to the limbic network being "sensory". An alternative explanation is that the limbic network essentially overlaps with regions of BOLD signal dropout, leading to low-SNR and flatter power spectrum (see e.g. SNR maps and discussion in (Yeo et al., 2011; Du et al., 2024)).}

We appreciate this insight. We have revised our discussion in \textbf{Section 5.2} to acknowledge that signal dropout and low SNR in the limbic network may contribute to the observed short timescales, citing Yeo et al. (2011) and Du et al. (2024).

\manuscriptquote{Finally, even with careful denoising, interpretational caution is warranted in regions susceptible to signal loss, such as the orbital frontal cortex and the anterior inferior medial temporal lobe (Thomas Yeo et al., 2011; Du et al., 2024). [Section 5.2]}

\reviewercomment{3. Motion censoring is an effective and commonly employed strategy in fMRI analysis, but its implementation in contexts where dynamics are of interest (as is the case for timescale analysis) is nontrivial; see e.g. (Goldberg et al., 2024). On the other hand, perhaps the AR1 projection method advocated by the authors would provide an elegant resolution to this matter by naturally providing a method for interpolation over censored frames? (Or perhaps censored frames can be excluded altogether during AR1 model fitting by following a procedure analogous to the "exact DMD" approach (see 2.2 in (Tu et al., 2014).)}

Thank you for this suggestion. We have highlighted in \textbf{Section 5.2} that our Time-Domain Linear Estimator allows for the exclusion of motion-contaminated frames without the need for interpolation, as it relies only on adjacent timepoint pairs.

\manuscriptquote{Regarding motion, the proposed Time-Domain Linear Estimator offers a practical advantage... this estimator relies solely on adjacent timepoint pairs $(x_t, x_{t+1})$. This allows motion-contaminated frames to be excluded without interpolation (Goldberg et al., 2024). [Section 5.2]}

\subsection*{Reviewer 2}

\reviewercomment{You treat timescales as intrinsic properties of the "aperiodic" component of the BOLD signal, yet this framing is both circular and conceptually fragile. In practice, 1/f slope, often used to operationalize "aperiodic", to the point of becoming a synonym, can emerge from entirely rhythmic processes with nonstationary amplitude envelopes. In this sense, the assumption that BOLD is mostly aperiodic is either difficult to verify (if we look for "truly" aperiodic activity, supposing it exists), or a truism, given that the amplitude of rhythmic processes will definitely change over several minutes.}

We agree that the distinction between ``aperiodic'' and ``rhythmic'' can be circular. We have clarified in \textbf{Section 1.2} that our method avoids external criteria to separate these components and is robust to mixed processes. We also added \textbf{Section 5.3 (Effects of Periodicity)} and \textbf{Figure 8} to demonstrate this robustness via simulations of AR2 processes with increasing oscillatory strength.

\manuscriptquote{Because amplitude-modulated rhythmic activity can mimic aperiodic properties (He et al., 2010; He, 2011), distinguishing the true generative process is often circular. Instead, we introduce a timescale estimator that is robust regardless of whether the underlying process is strictly aperiodic, damped-oscillatory, or a mixture of both... [Section 1.2]}

\reviewercomment{The slow BOLD fluctuations underlying the estimated timescales are highly susceptible to non-neuronal physiological influences, including systemic low-frequency oscillations (sLFOs) driven by autonomic processes such as respiration and CO₂ variation... Methods such as RAPIDtide... have shown that these delays are widespread and can dominate low-frequency variance in resting-state fMRI. Here you don't address this issue, which is particularly problematic given the centrality of slow autocorrelations in the analysis.}

We have addressed this by applying RAPIDtide to our HCP dataset to remove these physiological confounds, as described in \textbf{Section 4.1}.

\manuscriptquote{To further mitigate systemic low-frequency physiological oscillations driven by respiration and CO2 variation, we applied RAPIDtide (Frederick, 2025). This step accounts for the regionally varying lags of these global signals... [Section 4.1]}

\reviewercomment{A recent study by Korponay et al. ... shows that rsFC values inflate artifactually throughout fMRI scans with a substantial contribution from the mechanism explained above, and connected to arousal, beyond neural coupling. The same processes could inflate regional autocorrelation and thus timescale estimates. This questions whether the reported cortical hierarchies truly reflect neural temporal integration or are consequences of temporally correlated physiological fluctuations.}

We have cited Korponay et al. (2024) in \textbf{Section 5.3} and acknowledged that systemic oscillations can mimic neural timescales. We argue that while our estimator is statistically robust to these oscillations (as shown in our new simulations), valid biological interpretation requires effective denoising (like RAPIDtide).

\manuscriptquote{Furthermore, non-neural sources, such as respiration and CO2 variation, can introduce systemic low-frequency oscillations (Tong et al., 2019; Korponay et al., 2024). These dynamics manifest as periodic fluctuations in the ACF... [Section 5.3]}

\reviewercomment{The hemodynamic response function is another filter that displays spatial and subject specificity...}

We have added \textbf{Section 5.1 (Effect of the Hemodynamic Response Function)} to explicitly discuss how the HRF acts as a low-pass filter that inflates timescales. We note that while it preserves relative hierarchies, HRF variability can be a confound.

\manuscriptquote{The hemodynamic response function (HRF) acts as a temporal low-pass filter on this latent neural activity... reliance on a canonical HRF may oversimplify the relationship between neural and vascular dynamics. [Section 5.1]}

\reviewercomment{You correctly argue that standard errors derived from the ACF itself are biased under most realistic signal structures. However, this limitation is underplayed in the manuscript. The rationale for the hybrid autocorrelation/time-domain method needs to be better justified, particularly in emphasizing why modeling the error structure of ACFs is nontrivial and often invalid without explicit noise models...}

We have expanded the justification for the hybrid method in \textbf{Section 2.5.3}, explaining that ACF residuals in autoregressive processes represent misspecification error rather than random noise, necessitating time-domain variance estimation.

\manuscriptquote{For the special case of correct specification... the autocorrelation-domain residuals $e_k$ are zero, giving the false impression that the standard error vanishes... To address this, we propose a hybrid approach where the timescale is defined in the autocorrelation domain... but its standard error is defined in the time domain... [Section 2.5.3]}

\reviewercomment{Ultimately, what are we supposed to do with these fMRI timescales? If they should reflect neural timescales, I don't think we're there yet. If they are (yet another) derived property from BOLD signal, what does it add to the existing ones?}

We discuss the interpretation and utility of these maps in \textbf{Section 5}, emphasizing that while they align with functional hierarchies, mechanistic interpretation requires caution. We view these methods as a rigorous statistical foundation that enables future work to disentangle neural from vascular contributions.

\manuscriptquote{Nevertheless, a fundamental identifiability problem remains in that both a wider HRF and a longer neural timescale can produce nearly identical observed BOLD signals. Resolving this ambiguity requires moving beyond observational studies... [Section 5.1]}

\subsection*{Reviewer 3}

\reviewercomment{2.1 Distinguishing Measurement Notion from Innovation\\
The AR1 model at the heart of much of the work presented in (3) has a fundamental shortcoming: the "errors" $e_t$ are really innovations... In the fMRI setting, it seems far more plausible that the observation $X_t$ is really the sum of several distinct stochastic processes, at the very least one capturing neural dynamics, and another capturing measurement noise... I think the paper would benefit from more explicitly articulating this challenge, and perhaps exploring it in an additional simulation setting.}

We agree with this distinction. We have added \textbf{Section 5.2 (Effects of Measurement Noise)} and a new simulation study (\textbf{Figure 7}) modeled by Equations 54 and 55. This simulation explicitly separates latent neural dynamics (innovations) from added measurement noise. We show that measurement noise biases the estimated timescale downwards, and we discuss this as estimating a ``pseudo-true'' parameter.

\manuscriptquote{To illustrate this limitation, consider the simulation setting where the latent neural process is generated as an AR1... Observations are formed by adding independent measurement noise of increasing magnitude... This dependency emphasizes that the fitted timescale is effectively a confounded measure of both neural dynamics and measurement noise... [Section 5.2]}

\reviewercomment{3.1 Distinguishing Estimands from Estimators\\
The paper sometimes conceptually equivocates between estimands... and estimators...}

We have clarified this distinction throughout the text. In \textbf{Section 2.1.3}, we explicitly define the ``pseudo-true parameter'' (estimand) denoted by a superscript star, distinct from the estimator.

\manuscriptquote{Throughout, a superscript star (e.g. $\phi^*$) denotes the pseudo-true parameter corresponding to the best fitting parameter in a possibly misspecified model. [Section 2.1.3]}

\reviewercomment{3.2.1 Decouple Models from Estimation Techniques\\
The paper largely hinges on the "working model" that is used in analysis. However, the notation conflates this with the estimation technique employed...}

We have refined the notation to clearly distinguish between the model definitions and the estimation techniques. We retained ``TD'' and ``AD'' subscripts to denote the domains (Time-Domain vs Autocorrelation-Domain) which align with the working models, as suggested.

\reviewercomment{3.2.2 More Background Needed for Misspecification\\
While the paper gives a good amount of background on timescale estimation, it would be helpful to include a brief primer on estimation and inference under misspecification.}

We have added \textbf{Section 2.1.3 (Estimation and Inference under Misspecification)} to provide this background, defining the M-estimator approach and the minimization of squared error loss.

\reviewercomment{3.2.3 Inconsistent Use of $\cdot^*$\\
The authors frequently use $\cdot^*$ ... however, they are not entirely consistent...}

We have standardized the use of the $^*$ superscript to denote population-level pseudo-true parameters throughout the Methods section.

\reviewercomment{3.2.4 Xt vs xt\\
I’m a bit confused by what kind of objects Xt and xt are.}

We have clarified in \textbf{Section 2.1} that $X_t$ is the stochastic process and $x_t$ denotes an observed sample.

\manuscriptquote{Let $X_t$ ... be a discrete-time stochastic process... and let $x = \{x_1, x_2, . . . , x_T\}$ where $x_t$ denotes an observed sample of $X_t$. [Section 2.1]}

\reviewercomment{3.3 Preface Reliance on Hansen’s Textbook\\
The manuscript relies heavily on results from Hansen’s 2022 "Econometrics" textbook. While this is entirely fine, I think it would behoove the authors to preface and justify this reliance...}

We have added a preface to \textbf{Section 2 (Methods)} explicitly stating our reliance on Hansen (2022).

\manuscriptquote{There are varying definitions and conventions in the study of time series data... To simplify our presentation, this work closely follows the Hansen (2022) Econometrics textbook as a primary reference. [Section 2]}

\reviewercomment{3.4 Standard Errors Before Normality\\
I’d suggest rearranging the text so that much of §2.6 comes before §2.4 and §2.5. I think it would be more natural to first establish that the estimators are asymptotically normal...}

We have reordered the sections. \textbf{Estimator Properties (now Section 2.4)}, which establishes asymptotic normality, now precedes \textbf{Standard Error of the Estimators (Section 2.5)} and \textbf{Standard Error Estimation (Section 2.6)}.

\reviewercomment{4 Minor Issues}

Please see the revised manuscript where we have addressed all these minor concerns, including the clarification of temporal vs. spatial autocorrelation, notational fixes, and figure improvements.

\end{document}